\documentclass[11pt,a4paper]{article}
\usepackage[T1]{fontenc}
\usepackage[utf8]{inputenc}
\usepackage{amsmath,amssymb,xcolor}
\usepackage[margin=25mm]{geometry}
\usepackage[colorlinks=true,urlcolor=blue]{hyperref}
\definecolor{nexorablue}{HTML}{2459B5}
\setlength{\parindent}{0pt}
\setlength{\parskip}{7pt}
\setlength{\emergencystretch}{2em}
\renewcommand{\familydefault}{\sfdefault}
\begin{document}

{\small\bfseries\color{nexorablue} NEXORA CPGE}\par
{\small MPSI / NEX-LOG-03 / Corrigé}\par
{\LARGE\bfseries Un reste qui interdit une divisibilité}\par
\vspace{0.25em}\hrule\vspace{0.6em}

\section*{Contraposée}
Si $5$ divise $n$, alors $n^2-1\equiv-1\pmod5$, donc cette expression n’est pas divisible par $5$. C’est la contraposée recherchée.

\section*{Réciproque}
Elle affirmerait que tout entier non divisible par $5$ vérifie $5\mid(n^2-1)$. Pour $n=2$, on a $n^2-1=3$. La réciproque est fausse.


\end{document}
